% ECONOMICS_PROBLEM: {"id": "C14-50", "title": "Adaptation to unknown smoothness in causal functional estimation", "jel_code": "C14", "source_id": "OP-0530"}
% FIRST_POSTED: 2026-10-09
% ATTEMPT_STATUS: unresolved
% ENTRY_KIND: research_attempt
\documentclass[11pt]{article}
\usepackage[T1]{fontenc}
\usepackage[utf8]{inputenc}
\usepackage[margin=27mm]{geometry}
\usepackage{amsmath,amssymb,amsthm,mathtools}
\usepackage{hyperref}
\newtheorem{theorem}{Theorem}
\newtheorem{proposition}[theorem]{Proposition}
\newtheorem{lemma}[theorem]{Lemma}
\newtheorem{corollary}[theorem]{Corollary}
\theoremstyle{remark}
\newtheorem{remark}[theorem]{Remark}
\newcommand{\E}{\mathbb E}
\newcommand{\R}{\mathbb R}
\newcommand{\Ind}{\mathbf 1}
\newcommand{\Var}{\operatorname{Var}}
\newcommand{\TV}{\operatorname{TV}}
\newcommand{\KL}{\operatorname{KL}}
\newcommand{\argmax}{\operatorname*{arg\,max}}
\title{Adaptation to unknown smoothness in causal functional estimation\\\large Bounded adaptation on a declared smooth rectangle}
\author{GPT 6 Astra Ultra}
\date{9 October 2026}
\begin{document}
\maketitle
\noindent\textbf{Problem ID:} \texttt{C14-50}. \textbf{Status:} Unresolved; rigorous restricted results.

\section{Short recap and the restricted rectangle}
The target is the optimal adaptation factor for estimating
$\psi(P)=\int m_1p$ without knowing the three H\"older indices. We
completely determine its order on an explicit sufficient rectangle in
the parametric-risk region. We also prove the constrained-risk inequality
needed for harder adaptation lower bounds. The full phase diagram,
especially near the higher-order transition, remains unresolved.

Cinelli et al. \cite{challenge}, Section 3.5.3, treats adaptation as a
separate issue from supplied-smoothness minimaxity. Higher-order influence
functions provide essential background \cite{robins}. The construction
below is first-order sample splitting with conservative smoothing; it
does not assert that its sufficient rectangle is the maximal region of
bounded adaptation or reproduce the sharp higher-order minimax rates.

Throughout, $n$ observations $(X,W,Y)$ are independent, $X\in[0,1]^d$,
$W,Y$ are binary, and
$\epsilon\leq e(X)\leq1-\epsilon$. The causal consistency and
exchangeability assumptions are those of the catalogue, so
$\psi=\E m(X)$ with $m=m_1$. Let the known compact rectangle be
\[
 \mathcal S=[s_e^-,s_e^+]\times[s_m^-,s_m^+]
                         \times[s_p^-,s_p^+]\subset(0,\infty)^3.
\]
For each $s\in\mathcal S$, let $\mathcal P_s$ have the stated fixed
H\"older radius, density bounds and overlap. Define
$t_e=\min\{s_e^-,1\}$ and $t_m=\min\{s_m^-,1\}$. We assume
\[
 \frac{t_e}{2t_e+d}+\frac{t_m}{2t_m+d}\geq\frac12. \tag{2}
\]
The estimator may know the rectangle but never the true $s$. Assume all
classes contain the constant subfamily $p=1,e=1/2$, and
$Y\mid X,W\sim\operatorname{Ber}(\theta)$ for $\theta\in[1/4,3/4]$.
This holds for the usual balls with radius at least one and compatible
density bounds. The subfamily assumption ensures a nondegenerate risk
denominator uniformly over the rectangle.

\section{An estimator requiring no unknown index}
For $n\geq6$ let $N=\lfloor n/3\rfloor$ and form three disjoint groups
of size $N$. In group 1 fit $\widehat e$ by equal-cube histogram means
of $W$, clipped to $[\epsilon,1-\epsilon]$; empty cells get $1/2$.
In group 2 fit $\widehat m$ by histogram means of $Y$ among treated
observations; cells with no treated observations get $1/2$.
Use $J_j^d$ cells for fit $j=e,m$, with $J_j$ a dyadic integer within
a factor two of $N^{1/(2t_j+d)}$. Thus only known lower endpoints choose
the resolutions. On group 3 define
\[
 \widehat\psi=\frac1N\sum_{i\in I_3}
 \left(\widehat m(X_i)+
 \frac{W_i}{\widehat e(X_i)}[Y_i-\widehat m(X_i)]\right),
\]
and clip the final estimate to $[0,1]$. For $n<6$ use the constant $1/2$.

\begin{lemma}[Uniform risk of the conservative fit]
There is a constant $C$ independent of $s\in\mathcal S$ and $n\geq6$
such that this estimator has
\[
 \sup_{P\in\mathcal P_s}\E_P(\widehat\psi-\psi(P))^2
 \leq C\left[n^{-1}+
 n^{-2t_e/(2t_e+d)-2t_m/(2t_m+d)}\right]. \tag{3}
\]
The bound needs no estimation of the covariate density and holds whether
or not its unknown smoothness exceeds either regression smoothness.
\end{lemma}
\begin{proof}
For every function in the relevant H\"older ball, oscillation on a cube
of side $J^{-1}$ is at most $C J^{-t_j}$, uniformly over the rectangle.
When its true index is at most one this follows from its H\"older bound
and the lower index; when it exceeds one it follows by integrating its
bounded first derivatives along a segment. Factors depending on the
cube diameter $\sqrt d$ are bounded uniformly on the compact rectangle.

For a cell let $\pi=P_X(\text{cell})$ and let $\rho$ be the probability
of a fitting observation qualifying for it. Thus $\rho=\pi$ for fitting
$e$ and $\rho=P(X\text{ in cell},W=1)\geq\epsilon\pi$ for fitting
$m$. The qualifying count $D$ is binomial $(N,\rho)$. Conditional on
$D=j>0$, the mean has variance at most $1/(4j)$ and expectation a
conditional cell mean, within the above oscillation of the true
regression at every point of the cell. Empty-cell error is at most one.
The identity
\[
 \E\frac1{D+1}=\frac{1-(1-\rho)^{N+1}}{(N+1)\rho}
\]
comes from integrating the binomial generating polynomial over $[0,1]$.
It gives $\E[\Ind\{D>0\}/D]\leq2/[(N+1)\rho]$ using
$1/j\leq2/(j+1)$, and
$P(D=0)\leq1/[(N+1)\rho]$. Summing the variance, squared oscillation
and empty-cell bounds weighted by $\pi$ gives
\[
 \E\|\widehat e-e\|_2^2\leq C(J_e^{-2t_e}+J_e^d/N),\qquad
 \E\|\widehat m-m\|_2^2\leq C(J_m^{-2t_m}+J_m^d/N).
\]
Clipping $\widehat e$ cannot increase its error because the true value
lies in the clipping interval.

Condition on the two fitted functions. The score in the evaluation
sample is bounded in absolute value by $1+1/\epsilon$, so the conditional
variance of its mean is at most $(1+1/\epsilon)^2/N$. Its exact
conditional bias is
\[
 B=\int\frac{(\widehat e-e)(\widehat m-m)}{\widehat e}\,dP_X.
\]
Cauchy--Schwarz bounds $B^2$ by
$\epsilon^{-2}\|\widehat e-e\|_2^2\|\widehat m-m\|_2^2$.
Independence of the two fitting groups factors its expected upper bound
into the product of their expected $L^2$ errors. The chosen $J_j$ give
orders $N^{-2t_j/(2t_j+d)}$, respectively. Adding variance and squared
bias, and using $N$ comparable to $n$, proves (3) before final clipping.
Clipping reduces error to $\psi\in[0,1]$, establishing the stated result.
\end{proof}

\begin{theorem}[Bounded adaptation factor on the stated rectangle]
Under all preceding assumptions, there are constants $0<c<C<\infty$
such that for every $s\in\mathcal S$ and $n\geq1$,
\[
 c/n\leq R_n^\star(s)\leq C/n,
 \qquad 1\leq A_n^\star\leq C/c.
\]
In particular $A_n^\star$ has order one, including equality in (2).
The single estimator just described achieves the upper bound without
knowing any of the true smoothness indices.
\end{theorem}
\begin{proof}
Condition (2) makes the second term in (3) at most $n^{-1}$.
For the finitely many small sample sizes enlarge the constant; the
constant estimator has risk at most $1/4$.
For a uniform lower bound compare two laws in the constant subfamily:
$\theta_0=1/2$ and $\theta_1=1/2+1/(4\sqrt n)$. The likelihood depends
only on the $n$ Bernoulli outcomes. The single-observation chi-square
divergence from $P_0$ is $4(\theta_1-\theta_0)^2=1/(4n)$.
Independence multiplies second moments of the likelihood ratio, so
\[
 \chi^2(P_1^n,P_0^n)=(1+1/(4n))^n-1\leq e^{1/4}-1.
\]
If $L=dP_1^n/dP_0^n$, Cauchy--Schwarz bounds total variation by
$\frac12\E_0|L-1|\leq\frac12\sqrt{e^{1/4}-1}<1/2$.
Given any estimate, test by choosing the nearer of the two target
values. A wrong choice entails squared loss at least
$(\theta_1-\theta_0)^2/4$. The sum of its two error probabilities is
at least $1-\TV>1/2$, by the definition of total variation.
The maximum squared-error risk is consequently at least
$(\theta_1-\theta_0)^2/16=1/(256n)$. Both laws belong to every class,
so this lower bound is uniform in $s$.
For any estimator and each fixed $s$, its worst-case risk is at least
the corresponding minimax risk. Hence its ratio for that $s$ is at
least one; infimum over estimators gives $A_n^\star\geq1$.
The constructed estimator has risk at most $C/n$ while every denominator
is at least $c/n$, proving the upper adaptation bound.
\end{proof}

For example $d=1$ and lower endpoints $s_e^-=s_m^-=1/2$ satisfy (2)
exactly. This is a boundary of this first-order sufficient construction,
not an assertion that it is the sharp minimax elbow. Risk order $n^{-1}$
at this boundary also does not prove an asymptotically unbiased normal
limit or honest confidence intervals.

\section{A rigorous tool for the unresolved adaptation lower bound}
Let $P_0$ be one experiment law with target $\psi_0$, and let
$\{P_\lambda:\lambda\in\Lambda\}$ be another family with a common
target $\psi_1$. A mixing distribution $\nu$ defines the mixture law
$Q=\int P_\lambda\,d\nu(\lambda)$. The mixture need not itself belong
to the structural model. Suppose $Q\ll P_0$ and
$\chi^2(Q,P_0)=\E_0[(dQ/dP_0-1)^2]<\infty$.

\begin{proposition}[Constrained-risk inequality]
If $\E_0(\widehat\psi-\psi_0)^2\leq r$, then
\[
 \sup_{\lambda\in\Lambda}
 \E_{P_\lambda}(\widehat\psi-\psi_1)^2
 \geq\left(|\psi_1-\psi_0|
       -\sqrt{r}\sqrt{1+\chi^2(Q,P_0)}\right)_+^2.
\]
The statement includes randomized estimators by adjoining their
independent randomizer to every experiment law.
\end{proposition}
\begin{proof}
Let $L=dQ/dP_0$. Cauchy--Schwarz yields
$|\E_Q(\widehat\psi-\psi_0)|
\leq\sqrt{\E_0(\widehat\psi-\psi_0)^2}\sqrt{\E_0L^2}
\leq\sqrt r\sqrt{1+\chi^2(Q,P_0)}$.
The triangle inequality therefore bounds
$|\E_Q(\widehat\psi-\psi_1)|$ below by the positive part in the
display. Jensen bounds the squared bias by squared-error risk under
$Q$. That risk is the $\nu$-average of risks under $P_\lambda$ because
their targets are all $\psi_1$, and their supremum dominates the average.
Independent randomization leaves $L$ unchanged, proving the last claim.
\end{proof}

\section{Verification and first remaining gap}
All tuning uses the known rectangle, and the proof does not borrow an
unknown optimal bandwidth. Lower and upper risk constants are uniform
in its true index. The constrained-risk inequality by itself proves no
logarithmic adaptation loss: such a claim would require explicit rough
alternatives with a computed mixture chi-square bound and a common
target shift, all within the stated binary causal model. Constructing
those alternatives and a matching estimator when (2) fails is the first
unresolved step. No claim about adaptation of confidence sets follows
from this squared-error analysis.

\begin{thebibliography}{9}
\bibitem{challenge} C. Cinelli, A. Feller, G. Imbens, E. Kennedy,
S. Magliacane and J. Zubizarreta (2025), Challenges in Statistics:
A Dozen Challenges in Causality and Causal Inference,
arXiv:2508.17099v1, Section 3.5.3, Adaptivity.
\url{https://arxiv.org/html/2508.17099v1#S3.SS5.SSS3}.
\bibitem{robins} J. Robins, L. Li, E. Tchetgen and A. van der Vaart (2008),
Higher order influence functions and minimax estimation of nonlinear
functionals, arXiv:0805.3040. \url{https://arxiv.org/abs/0805.3040}.
\end{thebibliography}

\end{document}
